\chapter*{Remerciements}

Je tiens tout d'abord à remercier ISOGEO de m'avoir accueilli dans le cadre de
ma troisième année de formation d'ingénieur à Géodata Paris. Cette alternance
m'a permis d'évoluer dans un environnement professionnel consacré à la
valorisation, à la documentation et à la gestion des données géographiques,
tout en participant directement à l'évolution d'un produit utilisé dans des
contextes opérationnels variés.

J'adresse mes sincères remerciements à mon maître d'apprentissage,
\ReportCompanyTutor, pour la confiance qu'il m'a accordée, son accompagnement,
sa disponibilité et la qualité de ses conseils tout au long de cette année.
Ses retours m'ont permis de mieux comprendre les enjeux techniques et
fonctionnels du produit Scan, mais également d'améliorer progressivement ma
manière d'analyser, de concevoir et de valider une solution logicielle.

Je remercie également l'ensemble des membres de l'équipe Produit d'ISOGEO pour
leur accueil, leur disponibilité et les nombreux échanges qui ont accompagné
mon travail.
Leur connaissance du produit, des sources de données prises en
charge et des traitements historiques réalisés avec FME a été essentielle pour
comprendre l'existant et proposer des développements cohérents avec
l'architecture de Scan.

Je remercie tout particulièrement les personnes ayant participé aux échanges
techniques, aux revues de code et aux phases de validation.
Leurs remarques ont
contribué à améliorer la robustesse des traitements développés, la gestion des
cas particuliers ainsi que la qualité générale de la solution.

J'adresse également mes remerciements à mon professeur référent,
\ReportAcademicTutor, pour son suivi et ses conseils durant cette période
d'alternance.
Je remercie plus largement l'ensemble de l'équipe pédagogique de
Géodata Paris pour les connaissances et les méthodes transmises au cours de ma
formation d'ingénieur en géomatique.

Enfin, je remercie toutes les personnes qui m'ont soutenu et encouragé pendant
la réalisation de cette alternance et la rédaction de ce rapport.
